\documentclass[11pt]{article}
\usepackage[margin=1in]{geometry}
\usepackage[T1]{fontenc}
\usepackage{lmodern}
\usepackage{amsmath,amssymb}
\usepackage{hyperref}

\newcommand{\DPr}{\mathrm{DPr}}

\title{Literature Status for arXiv:math/0210287: the RNP Quotient Question}
\author{Run fa\_banach\_001, agent\_lane\_13}
\date{2026-06-16}

\begin{document}
\maketitle

\section*{Status}

This is a \texttt{literature\_already\_answered} packet.  It records that the
second final question in Kadets--Kalton--Werner, \emph{Remarks on rich
subspaces of Banach spaces}, arXiv:math/0210287, was later answered
negatively by Kadets--Shepelska--Werner, \emph{Quotients of Banach spaces with
the Daugavet property}, arXiv:0806.1815.  This packet also retains the earlier
local positive partition subcase as included history, not as a separate result
packet.

\section*{Source Question}

The source paper arXiv:math/0210287 ends with two questions of
A.~Pelczynski.  The second asks whether, if \(X\) is a subspace of \(L_1\) with
the Radon--Nikodym property, \(L_1/X\) must have the Daugavet property.

\section*{Supporting Answer}

Kadets--Shepelska--Werner arXiv:0806.1815 explicitly identifies this question
in its introduction, saying it was reiterated in Shvidkoy's paper and in
Kadets--Kalton--Werner.  The paper introduces poor subspaces and proves, in
Theorem~6.10 and the following corollary in the arXiv/published numbering, that
there is a subspace
\[
  E\subset L_1[0,1]
\]
which is isomorphic to \(\ell_1\), and therefore has the Radon--Nikodym
property, while
\[
  L_1[0,1]/E
\]
fails the Daugavet property.  This is exactly a negative answer to the source
question.

\section*{Included Partial History}

Before the supporting counterexample was identified, this run produced a
positive structured subcase, now stored inside this packet at
\[
\texttt{included\_partial\_history/0210287\_partition\_ell1\_rnp\_quotient\_daugavet/}.
\]
It proves that if \((A_n)\) is a countable measurable partition of a nonatomic
probability space and
\[
  X=\overline{\operatorname{span}}\{\mathbf 1_{A_n}/\lambda(A_n):n\in\mathbb N\}
  \subset L_1,
\]
then \(X\) is isometric to \(\ell_1\), hence has the Schur property and the
Radon--Nikodym property, and \(L_1/X\) has the Daugavet property.

This positive subcase is useful context but is subsumed by the later negative
literature answer.  The two results are compatible: partition-generated copies
of \(\ell_1\) can have Daugavet quotients, while arXiv:0806.1815 constructs a
different copy of \(\ell_1\) whose quotient fails the Daugavet property.

\section*{Scope}

This packet fully settles the second final question from arXiv:math/0210287 as
an already-known literature answer.  It does not address the separate first
question asking whether there exists a rich subspace of \(L_1\) with the Schur
property.

The included partition subcase should be read as provenance and contrast, not
as a separate registry result for the original question.

\section*{Search Evidence}

The decisive source was found by searching the local run indexes and web for
phrases including \texttt{poor subspace}, \texttt{L\_1/X Daugavet RNP},
\texttt{Radon-Nikodym subspace of L\_1 Daugavet property}, and the supporting
title.  The arXiv abstract of 0806.1815 states that a quotient of \(L_1[0,1]\)
over an \(\ell_1\)-subspace can fail the Daugavet property and that this answers
Pelczynski's question negatively.  The source text further names
arXiv:math/0210287 as one place where the question was reiterated.  The prior
local partial packet was moved inside this packet so the run has one canonical
package for the source question.

\begin{thebibliography}{9}
\bibitem{KKW}
V.~Kadets, N.~Kalton, and D.~Werner,
\emph{Remarks on rich subspaces of Banach spaces},
arXiv:math/0210287.

\bibitem{KSW08}
V.~Kadets, V.~Shepelska, and D.~Werner,
\emph{Quotients of Banach spaces with the Daugavet property},
arXiv:0806.1815; Bull. Pol. Acad. Sci. 56 no. 2 (2008), 131--147.

\bibitem{localpartial}
Run fa\_banach\_001, agent\_lane\_13,
\emph{A Partition \(\ell_1\) Subcase of the \(L_1/X\) Daugavet Question},
included in this packet under
\texttt{included\_partial\_history/0210287\_partition\_ell1\_rnp\_quotient\_daugavet/}.
\end{thebibliography}

\end{document}
